\subsection{Goldsheid's coordinate-product approximation: a counterexample and quantitative bounds}
\label{dev:1981-goldsheid}

\paragraph{The exact question and its two quantifiers.}
The intervention on P001018 (1981-557.8; GP-1981-40) concerns the standard
Euclidean symplectic space \(\mathbb R^{2n}\), with an orthonormal symplectic
basis \(e_1,\ldots,e_n,f_1,\ldots,f_n\). For every Lagrangian \(n\)-plane
\(L\), can one choose a nonzero vector in each prescribed plane
\(E_j=\operatorname{span}(e_j,f_j)\), so that their span \(K\) satisfies
\[
 \|P_L-P_K\|_{2\to2}\le 1/\sqrt2? \tag{G1}
\]
The notebook separately says that any constant \(c<1\), independent of
\(n\), would suffice. All projections here are Euclidean orthogonal
projections. Changing to a different, unquantified notion of distance would
change the problem.

\paragraph{Literature and historical limits.}
The searches in the accompanying log found no exact answer under the
speaker's name or the equivalent matrix formulation below. This is a
bounded negative retrieval result, not a proof that the problem has been
absent from the literature. A precise related result is
\citet[Theorem 3.4]{mehrmannpoloni2012doubling}: after symplectic coordinate
permutations, a Lagrangian has a graph matrix whose diagonal entries have
absolute value at most one and whose off-diagonal entries have absolute
value at most \(\sqrt2\). The theorem controls entries, and hence gives
dimension-dependent norm estimates; it does not state (G1) or a uniform
operator-gap estimate. We found no source explaining why Goldsheid's
specific question was dropped or whether it was later reformulated.

The results below are proved in this task, without a claim of historical
priority. They disprove (G1), show that any universal constant must be at
least \(\cos(\pi/8)\), and give an upper bound for every fixed dimension.
The still plausible weaker formulation is the existence of some
\(c\in[\cos(\pi/8),1)\) working in every dimension; no particular value
in that interval is asserted to work.

\begin{lemma}[Symmetric-unitary formulation]\ours
Identify \(\mathbb R^{2n}=\mathbb C^n\), with \(f_j=i e_j\), and let
\(Cz=\overline z\). Every real Lagrangian \(L\) is \(U\mathbb R^n\)
for a unitary matrix \(U\). Set \(S=UU^T\). Then \(S\) is symmetric
unitary and, as real-linear operators,
\[
 P_L=\tfrac12(I+SC).
\]
For \(K=D_\theta\mathbb R^n\), where
\(D_\theta=\operatorname{diag}(e^{i\theta_j})\),
\[
 \|P_L-P_K\|=\tfrac12\|S-D_\theta^2\|. \tag{G2}
\]
Consequently (G1) is equivalent to choosing a diagonal unitary \(D\)
such that \(\operatorname{Re}(S^*D)\) is positive semidefinite, where
\(\operatorname{Re}B=(B+B^*)/2\).
\end{lemma}
\begin{proof}
Choose a real orthonormal basis of \(L\). Its Hermitian Gram matrix is
the identity: the real part is the Euclidean Gram matrix and the imaginary
part vanishes because \(L\) is isotropic. Thus its columns form a unitary
\(U\). Reflection across \(L\) is the real-linear involution
\(UCU^*=(UU^T)C\). Averaging it with the identity gives the projection.
The same formula for \(D_\theta\) gives (G2), because conjugation is an
isometry and the norm of a complex-linear matrix is the same whether
regarded as a real or complex operator. Every diagonal unitary has a
diagonal unitary square root, so no choices are lost. Finally
\[
 (S-D)^*(S-D)=2I-S^*D-D^*S
             =2I-2\operatorname{Re}(S^*D),
\]
which proves the last equivalence.
\end{proof}

\begin{theorem}[An explicit six-plane refutes the proposed constant]\ours
Let \(\omega=e^{2\pi i/3}\) and
\[
 V=\frac1{\sqrt3}
 \begin{pmatrix}1&1&1\\1&\omega&\omega^2\\1&\omega^2&\omega\end{pmatrix}.
\]
The real six-dimensional subspace
\[
 L=\{(V\overline z,z):z\in\mathbb C^3\}\subset\mathbb C^6
       =\mathbb R^{12}
\]
is Lagrangian, and every coordinate-product Lagrangian \(K\) satisfies
\(\|P_L-P_K\|>1/\sqrt2\).
\end{theorem}
\begin{proof}
The Fourier matrix \(V\) is unitary. For \(z,w\in\mathbb C^3\), the
Hermitian inner product of the corresponding vectors in \(L\) is
\(\overline{z^*w}+z^*w\), which is real. The parametrization is real
linear and injective, so \(L\) is isotropic of half the ambient real
dimension and hence Lagrangian. An explicit real orthonormal basis is
\[
 \frac{(Ve_j,e_j)}{\sqrt2},\qquad
 \frac{(-iVe_j,ie_j)}{\sqrt2},\qquad j=1,2,3.
\]
Its reflection matrix is
\[
 S=\begin{pmatrix}0&V\\V^T&0\end{pmatrix}.
\]
Indeed \(S\) is symmetric unitary, and \(SC\) fixes exactly the displayed
plane: \(V^T\overline V=I\) supplies the second block equation.

Suppose a diagonal unitary \(D=\operatorname{diag}(d_1,\ldots,d_6)\)
made \(\|S-D\|\le\sqrt2\). The preceding lemma says
\(H=(S^*D+D^*S)/2\ge0\). Since every diagonal entry of \(S\) is zero,
every diagonal entry of \(H\) is zero. A positive semidefinite Hermitian
matrix obeys \(|H_{ij}|^2\le H_{ii}H_{jj}\), so \(H=0\).
For a nonzero entry \(S_{ij}\), this implies
\[
 d_i d_j=-\frac{S_{ij}}{\overline{S_{ij}}}. \tag{G3}
\]
The first row and first column of \(V\) are positive real. Their
equations force \(d_1=d_2=d_3=a\), \(d_4=d_5=d_6=b\), and \(ab=-1\).
The entry \(V_{22}=\omega/\sqrt3\) instead forces
\(ab=-\omega^2\), a contradiction. Thus the inequality cannot hold.
Compactness of the diagonal-unitary torus also shows that the minimum
gap for this one plane is strictly larger than \(1/\sqrt2\), not merely
that the infimum fails to be attained.
\end{proof}

\begin{theorem}[A Fourier-family obstruction to universal constants]\ours
For every odd integer \(m\ge3\), put
\(V_{jk}=m^{-1/2}e^{2\pi i jk/m}\), \(0\le j,k<m\), and form the
Lagrangian \(L_m=\{(V\overline z,z):z\in\mathbb C^m\}\subset\mathbb C^{2m}\).
Every coordinate-product Lagrangian \(K\) satisfies
\[
 \|P_{L_m}-P_K\|\ \geq\
 \frac12\sqrt{\,2+\sqrt{2-2/\sqrt m}\,}. \tag{G4}
\]
In particular any dimension-independent constant in Goldsheid's weaker
question must satisfy \(c\ge\cos(\pi/8)\).
\end{theorem}
\begin{proof}
As above, \(S=\bigl(\begin{smallmatrix}0&V\\V^T&0\end{smallmatrix}\bigr)\).
Write a diagonal unitary as \(D=\operatorname{diag}(A,B)\), with diagonal
entries \(a_j,b_k\). Then
\[
 \operatorname{Re}(S^*D)=\frac12
 \begin{pmatrix}0&C_0\\C_0^*&0\end{pmatrix},\qquad
 C_0=\overline V B+A^*V.
\]
The eigenvalues of the displayed off-diagonal Hermitian matrix are the
positive and negative singular values of \(C_0\), divided by two.
It follows exactly, without an estimate at this step, that
\[
 \|S-D\|^2=2+\|C_0\|. \tag{G5}
\]
Let \(W=(V_{jk}^2)\), the entrywise square, not the matrix square.
Because \(m\) is odd, multiplication by two permutes the residue classes
modulo \(m\). Consequently \(W\) is \(m^{-1/2}\) times a unitary
Fourier matrix with permuted columns, so \(\|W\|=m^{-1/2}\).
Expanding the Frobenius norm gives
\[
 \|C_0\|_F^2
 =2m+2\operatorname{Re}\sum_{j,k}\overline{a_j}\overline{b_k}V_{jk}^2
 \ge2m-2\sqrt m,
\]
because \(\|a\|=\|b\|=\sqrt m\) and
\(|a^*W\overline b|\le\|a\|\|W\|\|b\|=\sqrt m\).
There are \(m\) singular values, so
\(\|C_0\|\ge\|C_0\|_F/\sqrt m\ge\sqrt{2-2/\sqrt m}\).
Equations (G2) and (G5) prove (G4). Its right side tends to
\(\tfrac12\sqrt{2+\sqrt2}=\cos(\pi/8)\) as odd \(m\to\infty\).
Every \(c\) strictly below that limit is therefore defeated by some
finite odd \(m\). The endpoint itself is not excluded by this estimate.
\end{proof}

\begin{proposition}[A fixed-dimension upper bound]\ours
Every Lagrangian \(L\subset\mathbb C^n\) admits a common-phase product
plane \(K=e^{i\theta}\mathbb R^n\) with
\[
 \|P_L-P_K\|\le\cos\!\left(\frac{\pi}{2n}\right). \tag{G6}
\]
In particular the original bound is true for \(n=1,2\).
\end{proposition}
\begin{proof}
The eigenvalues of the unitary \(S\) lie on the unit circle. Of the \(n\)
circular gaps, including the wraparound gap and allowing repetitions,
one has length at least \(2\pi/n\). All eigenvalues therefore lie in a
closed complementary arc of length at most \(2\pi-2\pi/n\).
Let \(z=e^{i\phi}\) be its midpoint. Their angular distances from \(z\)
are at most \(\pi-\pi/n\), so, by the spectral theorem,
\[
 \|S-zI\|=\max_{\lambda\in\operatorname{Spec}S}|\lambda-z|
 \le2\sin\!\left(\frac\pi2-\frac{\pi}{2n}\right).
\]
Choose \(\theta=\phi/2\) and use (G2). If \(n=1\), the complementary
arc is a point and the bound is zero, as required.
\end{proof}

\begin{example}[Sharpness in dimension two]\ours
For \(S=\bigl(\begin{smallmatrix}0&1\\1&0\end{smallmatrix}\bigr)\),
each column of \(S-D\) has norm \(\sqrt2\) for every diagonal unitary
\(D\), so every product plane has gap at least \(1/\sqrt2\).
Taking \(D=iI\) gives equality, since \(S\) is self-adjoint unitary.
The associated Lagrangian is spanned over \(\mathbb R\) by
\((e_1+e_2)/\sqrt2\) and \((f_1-f_2)/\sqrt2\).
Thus (G6) is sharp at \(n=2\), although (G1) fails in higher dimension.
\end{example}

\begin{proposition}[The notebook's coordinatewise observation]\ours
For every coordinate plane \(E_j\), there is a unit vector
\(\varphi_j\in E_j\) with \(\operatorname{dist}(\varphi_j,L)\le1/\sqrt2\).
This does not bound the gap of the span of the chosen vectors.
\end{proposition}
\begin{proof}
Since \(L\) is Lagrangian, \(JL=L^\perp\), hence
\(P_L+JP_LJ^{-1}=I\). Compressing to the \(J\)-invariant plane \(E_j\)
and taking traces shows that the compression of \(P_L\) has trace one.
Its largest eigenvalue is at least \(1/2\), which proves the first
assertion by choosing a unit eigenvector.

For the second assertion, take \(n\ge2\),
\(v=(1,\ldots,1)/\sqrt n\in\mathbb R^n\), and
\(L=\mathbb Rv\oplus i v^\perp\). Its Hermitian pairings have zero
imaginary part and its real dimension is \(n\), so it is Lagrangian.
Every \(f_j=ie_j\) is at distance \(1/\sqrt n\) from \(L\), but
\(iv\in i\mathbb R^n\cap L^\perp\). Thus the span
\(K=i\mathbb R^n\) has \(\|P_L-P_K\|=1\). The equality follows also
from the fact that differences of orthogonal projections have norm at
most one. This example refutes an inference from a particular selection
of vectors, whereas the Fourier example refutes the existential (G1).
\end{proof}

\begin{remark}[Independent checks and numerical scope]\oursComputed
The independent mathematical audit re-derived the antiunitary reflection,
the positive-semidefinite zero-diagonal argument, the Fourier cross-term
estimate, the arc bound including \(n=1\), and the coordinatewise example.
It confirmed every hypothesis and emphasized that the endpoint
\(\cos(\pi/8)\) is not ruled out. A second representation of each Fourier
plane by the explicit real orthonormal basis checked isotropy directly.

The included \texttt{verify\_root.py} uses two independent random seeds,
five dimensions and four samples per dimension to compare the real
projection matrices with (G2); the maximum discrepancy was below
\(7\times10^{-16}\). It also checks Fourier examples for
\(m=3,5,7,11\). The rigorous lower bound for \(m=3\) is approximately
\(0.854312835963\), already above \(1/\sqrt2\). Two exploratory
optimizations found feasible gaps near \(0.866026\); neither met its
convergence criterion, and neither is reported as a global optimum.
The proofs of (G4)--(G6) do not depend on floating-point evidence. The
earlier random exploration, whose sampled gaps were all below
\(1/\sqrt2\), is also included: it illustrates why that evidence could
not justify the proposed universal bound.
\end{remark}

\paragraph{What remains and a bounded next investigation.}
Define
\[
 c_*:=\sup_{n\ge1}\ \sup_{S=S^T,\ S^*S=I}
       \inf_{D\ \mathrm{diagonal\ unitary}}\frac12\|S-D\|.
\]
The proofs establish \(\cos(\pi/8)\le c_*\le1\), with a strictly
subunit bound in each fixed dimension. Goldsheid's weaker question is
exactly whether \(c_*<1\). The original stronger assertion
\(c_*\le1/\sqrt2\) is false. A further bounded research run could
determine the exact optimum of the explicit \(m=3\) Fourier example by
an algebraic or interval certificate, or improve the Fourier-family
lower bound using higher singular-value moments. Either target is
concrete. A numerical optimizer alone would not decide the supremum over
all dimensions, and no mechanism for a uniform upper bound has been
proved here.

\subsection{Verification annotations for the spectral formulas}
\label{dev:1981-spectral-checks}

This subsection verifies corrections to established spectral theory. It
does not promote the Wodzicki exposition to a newly open C/D problem.

\begin{proposition}[Correct coefficients in the shifted-zeta identity]\ours
Let \(A\) be an invertible first-order elliptic differential operator on a
closed \(n\)-manifold, with a fixed admissible spectral cut. For a scalar
\(t\) sufficiently small that the cut and invertibility persist, put
\(\delta_\mu=\operatorname*{res}_{s=\mu}\zeta_A(s)\). For \(k\ge0\),
\[
\begin{split}
 \zeta_{A+t}(-k)={}&\sum_{\nu=0}^{k}\binom{k}{\nu}t^\nu\zeta_A(\nu-k)\\
 &+\sum_{\mu=1}^{n}
 \frac{(-1)^\mu\delta_\mu}{\mu\binom{k+\mu}{k}}\,t^{k+\mu}.
\end{split}
\tag{Z1}
\]
The second sum in P001025 (20) omits \(t^\mu\).
\end{proposition}
\begin{proof}
For sufficiently small \(t\) the complex-power binomial expansion,
initially in a trace-class half-plane, is
\[
 \zeta_{A+t}(s)=\sum_{j=0}^\infty
       \frac{(-1)^j(s)_j}{j!}t^j\zeta_A(s+j),
 \qquad (s)_j=s(s+1)\cdots(s+j-1).
\]
One may alternatively take a finite expansion with a remainder of
arbitrarily negative operator order, which is holomorphic near the point
under consideration. For \(0\le j\le k\), the coefficient at
\(s=-k\) is \(\binom{k}{j}\). For \(j=k+\mu\), \(\mu\ge1\),
the Pochhammer product has a simple zero and
\[
 \left.\frac{d}{ds}\frac{(-1)^{k+\mu}(s)_{k+\mu}}{(k+\mu)!}
 \right|_{s=-k}
 =\frac{(-1)^\mu k!(\mu-1)!}{(k+\mu)!}
 =\frac{(-1)^\mu}{\mu\binom{k+\mu}{k}}.
\]
It contributes only when \(\zeta_A(s+j)\) has a pole at \(\mu\).
For a differential operator of order one the positive integer poles are
confined to \(1,\ldots,n\), and the nonpositive integer values are
regular. Multiplying the zero by the pole gives (Z1). The factor \(t^j\)
has remained \(t^{k+\mu}\) throughout. The holomorphic remainder makes
no additional contribution, since its binomial coefficient vanishes at
\(-k\) once the expansion is taken beyond \(k+n\).
\end{proof}

\begin{remark}[Second check and the limit of the inference]\oursComputed
Direct differentiation of the finite Pochhammer product agrees with the
displayed coefficient in all 42 exact rational checks
\(0\le k\le6\), \(1\le\mu\le6\). Setting \(t=0\) independently
detects the notebook's missing powers, whenever the residue combination
is nonzero. Kernel-crossing cases are excluded from the proposition;
adding \(h(A)\) at zero without tracking which zero modes move is not a
substitute for that hypothesis.

The check also prevents an invalid inference about spectral cuts.
Cut independence at \(s=0\) for every shifted operator does not imply
cut independence at every negative integer. In the difference of (Z1)
for two cuts, the case \(k=0\) contains no negative values from which to
start such an induction. The general valid identity is instead
\citet[Theorem 4.2]{ponge2006spectral},
\(\rho(k)=(2\pi i/m)\operatorname{Res}(\Pi P^{-k})\).
Moreover, negative poles of eta and absolute zeta functions occur for
Dirac-type differential operators in odd dimensions
\citep[\S4]{loyamoroianu2012singularities}. For self-adjoint \(P\), the
two cut functions have the form
\(\zeta_+(s)+e^{\pm i\pi s}\zeta_-(s)\); their difference at an
integer is the limit of \(2i\sin(\pi s)\zeta_-(s)\).
A nonzero residue of \(\zeta_-\) therefore gives a nonzero difference,
even though each cut function is regular there. This is why the blanket
claim in P001025(d) is recorded as false without further hypotheses.
\end{remark}

\begin{theorem}[An explicit first-order differential operator with cut-dependent negative value]\ours
\label{dev:1981-negative-cut}
Let \(\mathbb T^3=(\mathbb R/2\pi\mathbb Z)^3\) have its flat metric,
let \(\mu>0\), and let \(\sigma_1,\sigma_2,\sigma_3\) be the Pauli
matrices. On the trivial Hermitian bundle of rank four set
\[
 \alpha_j=\sigma_1\otimes\sigma_j\quad(1\le j\le3),
 \qquad \beta=\sigma_3\otimes I_2,
 \qquad
 P=\sum_{j=1}^3\alpha_j(-i\partial_{x_j})+\mu\beta.
\]
Then \(P\) is an invertible self-adjoint elliptic differential operator
of order one and of Dirac type. Define \(Z_+(s)\) using argument
\(+\pi\) for negative eigenvalues, and \(Z_-(s)\) using argument
\(-\pi\); both branches give argument zero to positive eigenvalues.
The respective branch ranges are \((-\pi/2,3\pi/2)\) and
\((-3\pi/2,\pi/2)\), so the cuts are the negative and positive
imaginary axes. Both cut zeta functions are regular at \(s=-1\), but
\[
 \boxed{Z_+(-1)=-i\pi^2\mu^4,\qquad
        Z_-(-1)=+i\pi^2\mu^4.} \tag{Z2}
\]
In particular, their difference there is \(-2i\pi^2\mu^4\ne0\).
\end{theorem}
\begin{proof}
The four matrices \(\alpha_1,\alpha_2,\alpha_3,\beta\) are
Hermitian, square to the identity, and anticommute in distinct pairs.
Therefore
\[
 P^2=(-\Delta+\mu^2)I_4.
\]
The principal symbol squares to \(|\xi|^2I_4\), proving ellipticity
and the Dirac-type assertion. On the Fourier mode indexed by
\(k\in\mathbb Z^3\), the Hermitian matrix
\(\sum_j\alpha_jk_j+\mu\beta\) has square
\((|k|^2+\mu^2)I_4\) and trace zero. Its eigenvalues are
\(\pm\sqrt{|k|^2+\mu^2}\), each with multiplicity two. Fourier
expansion supplies the self-adjoint realization with domain \(H^1\),
and \(\mu>0\) excludes a kernel. Both stated cuts are admissible for
the operator and its principal symbol.

For \(\operatorname{Re}s>3\), put
\[
 Z_{\mathrm{abs}}(s)=\operatorname{Tr}|P|^{-s}
 =4\sum_{k\in\mathbb Z^3}(|k|^2+\mu^2)^{-s/2}.
\]
The multiplicities and branch conventions give exactly
\[
 Z_\pm(s)=\frac{1+e^{\mp i\pi s}}2\,Z_{\mathrm{abs}}(s). \tag{Z3}
\]
These identities continue meromorphically. To determine the pole of
\(Z_{\mathrm{abs}}\), use its Mellin representation
\[
 Z_{\mathrm{abs}}(s)
 =\frac1{\Gamma(s/2)}\int_0^\infty
 t^{s/2-1}\operatorname{Tr}(e^{-tP^2})\,dt.
\]
Gaussian Poisson summation on this torus yields
\[
\begin{split}
 \operatorname{Tr}(e^{-tP^2})
 &=4e^{-\mu^2t}\sum_{k\in\mathbb Z^3}e^{-t|k|^2}\\
 &=4\pi^{3/2}t^{-3/2}
 \left(1-\mu^2t+\frac{\mu^4t^2}{2}+O(t^3)\right)
   +O\!\left(t^{-3/2}e^{-c/t}\right)
\end{split}
\]
for some \(c>0\) as \(t\downarrow0\). The coefficient of
\(t^{1/2}\) is \(a=2\pi^{3/2}\mu^4\). The corresponding Mellin
integral contributes \(2a/(s+1)\) before division by
\(\Gamma(s/2)\). Since \(\Gamma(-1/2)=-2\sqrt\pi\),
\[
 \operatorname*{res}_{s=-1}Z_{\mathrm{abs}}(s)
 =\frac{2a}{\Gamma(-1/2)}=-2\pi\mu^4. \tag{Z4}
\]
After the displayed heat terms are subtracted, the small-time integral
is holomorphic near \(-1\); the large-time integral is entire because
\(\mu>0\). Thus the pole computation is justified without a formal
exchange of divergent series.

Each prefactor in (Z3) has a simple zero at \(-1\), with
\[
 \left.\frac{d}{ds}\frac{1+e^{\mp i\pi s}}2\right|_{s=-1}
 =\pm\frac{i\pi}{2}.
\]
Multiplying this derivative by the residue in (Z4) proves regularity
and the two values (Z2). The holomorphic part of
\(Z_{\mathrm{abs}}\) contributes zero at that point.
\end{proof}

\begin{remark}[Independent algebraic and normalization checks]\oursComputed
The separate executable \texttt{verify\_negative\_cut.py} checks all
Clifford anticommutators, Hermitian symmetry, the square and the exact
characteristic polynomial, the Mellin residue, and both branch limits.
The heat normalization was independently checked as
\(4\operatorname{vol}(\mathbb T^3)/(4\pi)^{3/2}=4\pi^{3/2}\).
A second mathematical audit reproduced the residue and phase signs.
The symmetric spectrum also gives \(\eta(P;s)=0\) identically:
this example uses the nonzero residue of the absolute zeta function.
It is an explicit verification of the notebook correction, with no
historical priority claim.
\end{remark}

\begin{remark}[The vector-calculus sign check]\oursComputed
For conventional Euclidean operators,
\(\operatorname{div}\operatorname{curl}=0\),
\(\operatorname{curl}\operatorname{grad}=0\), and
\(\operatorname{curl}\operatorname{curl}
=\operatorname{grad}\operatorname{div}-\Delta\).
Direct multiplication gives the corrected square in 1981-558.6 and
1981-559.12. An independent symbolic polynomial calculation replaces
the derivatives by three commuting indeterminates and reproduces both
the erroneous matrix's actual square and the corrected square exactly.
The output is included with the coefficient checks above.
\end{remark}
